\begin{annexitem}{annex:p12.en.state-invariant}{thm:p12-en-state-invariant}{Solved domains, the exact accumulator update and retained target scope}
\begin{proof}
Initially S and C are empty. Suppose the invariant holds at a
reachable state. Trivial, Orient and Decomposition modify only C
and introduce no variable not already there. Constrain inserts an
A-instance of an old equation. Since A's range contains no member
of its domain, this new equation contains none of \(D_S\).
At LP the fresh identity variables are outside \(D_S\). Its
matching constraint is \(l A=^{?}u A\), also without \(D_S\).
Thus these operations preserve the claimed constraint-domain property.

At elimination x occurs in C and hence is outside \(D_S\).
Its term t is a subterm of C, so contains no member of \(D_S\),
and the occurs proviso excludes x. In \(S\delta\) the old
left variables remain distinct and unchanged, while every old
right side loses x. The added binding has the new left variable x,
and t contains neither old nor new left variables. S' is solved
on \(D_S\cup\{x\}\). The updated C loses x and contains none
of that new solved domain.

For an old solved left variable d, both the new substitution and
\(A\delta\) give \((d A)\delta\). For x, A fixes x and the
composite gives t, the new binding's right side. Outside the new
domain both maps are identity. Hence \(\sigma_{S'}=A\delta\)
globally. The argument is identical for B and N; whether P receives
\(\delta\) does not alter these literal binding calculations.
Induction proves the reachable invariant and idempotence.
For N, initially P has no solved domain variables. Its constraint
rules leave P unchanged except at elimination, which removes x
and inserts t with no old solved-domain variable. Constrain only
removes an equation, and LP replaces a term by a fresh r with no
old solved-domain variable. Thus the new solved domain remains
absent from P at every N step. Applying A to P has no effect.

If \(\rho\) literally solves S, then for each d in \(D_S\)
its equation gives \((d A)\rho=d\rho\). Outside \(D_S\),
A is identity. Therefore \(A\rho=\rho\) on the whole variable
alphabet. All substitutions have finite support, so \(\rho\)
is a legitimate finite factor. In particular, when deleting a
literal identity removes an initial variable from the current
state, the same global calculation still covers its externally
retained target value. No claim that every source transformation
preserves the set of occurring variables is needed. For example,
Constrain and Trivial can delete the sole problem x=x; identity
still covers an arbitrary x-image by that image's factor.
\end{proof}
\end{annexitem}

\begin{annexitem}{annex:p12.en.trace-transport}{thm:p12-en-trace-transport}{Every primitive transition and the matching obligation commute with instantiation}
\begin{proof}
Induct on the finite basic trace. Initially A is identity, giving
identical P,C,S. Assume corresponding states satisfy
\(P_N=P_B A\) and have the same C,S. Occurrences in P are
retained with multiplicity; two equal instantiated equations do
not destroy the identification of the selected basic occurrence.

Trivial, Orient and Decomposition act on the same literal C
equation in both states. They have the same root/variable conditions,
produce identical C and leave P,S unchanged, so the relation persists.
At elimination use the same \(x=^{?}t\), the same \(\delta\),
and the same S'. B leaves \(P_B\) alone, while N produces
\(P_N\delta=P_B A\delta=P_B\sigma_{S'}\), by the preceding
global accumulator calculation. C and S' are identical. This proves
the exact elimination case; it would be wrong to leave \(P_N\)
unchanged or to compose in the opposite order.

For Constrain, let e be the chosen basic equation and \(e A\)
its corresponding ordinary equation. N inserts
\((e A)A=e A\), because A is idempotent; B inserts exactly
\(e A\). Removing the selected equation from each occurrence
multiset commutes with applying A. Thus both C,S and the remaining
P relation agree.
The optional bare-equation N Constrain also inserts exactly e A,
because the selected ordinary equation is fixed by A. More generally
the preceding N invariant proves equivalence of the two Constrain
variants even at N states not obtained by compiling a basic trace.

For LP choose the same fresh \(l\approx r\in\mathcal L\)
and the same basic nonvariable address p. Substitution cannot
remove a nonvariable root or change its arity, so at p the ordinary
term has subterm \(u A\), with exactly u's root. The root proviso
for nonvariable l is therefore the same in both systems. Freshness
gives \(l A=l\) and \(r A=r\). B inserts the matching constraint
\(l A=^{?}u A\); N inserts
\(l A=^{?}(u A)A\), which is literally the same constraint.
Substitution in the unchanged context commutes with hole replacement:
\[
                         (e[r])A=(e A)[r].
\]
This is exactly the equality of the new basic instantiated equation
and the actual ordinary replacement equation. The unselected equations
still agree after A, and S is unchanged. If l is a variable no
matching-first Decomposition is required in either representation.
If l is nonvariable, both sides of the same new constraint have
matching roots. Record the same matching-first obligation in both
traces. If the given basic trace first performs an operation elsewhere,
transport it by the corresponding case without moving Decomposition
earlier. When that trace first acts on the matching constraint, its
legal step is Decomposition; perform it on the same constraint in N.
Its argument constraints and subsequent obligation status are identical.
Thus no extra step is inserted into an arbitrary given basic trace.
For composite witness traces that deliberately decompose immediately,
the same immediate step is transported. An identical constraint already
present in the set C needs only one such Decomposition; the same
set and obligation occur in both representations. This does not
permit LP into an instantiated image that had no basic nonvariable
address.

The ordinary trace introduces no variable not introduced by the
basic trace; variables of instantiated terms already occur in the
basic state or its solved bindings. Consequently each variant fresh
for the earlier basic computation is also fresh for the transported
ordinary computation. Constraint Clash/Occurs, if the explicit
failure rules are included, also transport since C is identical.
This covers every primitive source case. Iterating the induction
constructs the finite N trace, not merely its answer.

At a successful terminal basic state P=C=empty, the corresponding
ordinary P is also empty and its C,S are identical. The returned
substitutions are therefore identical, including all fresh solved
bindings. Any factor of the basic answer on the initial variable
scope is literally the same factor of the ordinary answer. Thus the
source's passage from a successful B witness trace to Corollary4.14
has no additional normalization or convergence dependency. This
does not prove an opposite trace transport: a new position inside
an A-image may be available in N but not in the retained P_B.

In the new inverse-branch example, the first fresh identity is
the reverse of the second R-rule, with nonvariable left root g.
Both root and address restrictions hold at g(x). Its constraint
g(v)=g(x) decomposes first, then the binding v=x is solved.
The retained basic term is f(v,v); instantiation by S makes
it f(x,x), exactly as required for the next fresh forward
f(p,q)=a match. Decomposing f(p,q)=f(x,x) gives p=x,q=x,
and both eliminations pass occurs checks and keep solved form.
After Solve/Trivial on a=a, only fresh domains v,p,q remain
bound. The initial x is fixed, proving the claimed identity
answer in raw inverse-basic B and in its transported N trace.
The smaller forward family lacks that first g(v)=f(v,v)
identity, so the argument does not erase the directed failure.
\end{proof}
\end{annexitem}

\begin{annexitem}{annex:p12.en.mgu-blocks}{thm:p12-en-mgu-blocks}{Concrete helper runs and their exact literal factor conditions}
\begin{proof}
The reachable invariant ensures that Q contains no \(D_S\).
Run the actual source U helper on Q. Its variables remain among
\(\operatorname{Var}(Q)\); solved form and termination were proved
for that helper in \cref{thm:p12-eu-u-bridge}. Hence its output
\(\eta\) has domain and range variables outside \(D_S\), fixes
all other variables and is idempotent. For a Narrow match with
nonvariable l, roots agree by the LP proviso. Begin with Decomposition.
This is a legal U step and preserves the solution set; even an
identical nonvariable match can start with that step rather than
Trivial. Thus the helper respects the source's matching-first
restriction, not an unconstrained MGU oracle.

To realize this helper inside the source constraint state, follow
its operations on Q while applying each elimination substitution
also to old C and the right sides of old S. Old solved domains
are never selected, because Q has none. Each helper step remains
enabled in the surrounding state and updates the original bindings
by the same substitution. At completion the old context is
\(C\eta\) and the old bindings are \(S\eta\); the new bindings
are \([\eta]\). If a selected new constraint duplicates an old
set element, one concrete deletion may also remove that old
element. The removed element is solved by \(\eta\), so its
image in the printed \(C\eta\) is a literal identity. Deleting
such identities by Trivial reconciles the two outputs. No other
old constraint is lost. With C empty, the end states agree
literally and have C empty, precisely the source's composite
block boundaries.

Since \(\eta\) has no variable of \(D_S\), applying it to old
right sides keeps old solved left variables distinct. Its range
contains neither old solved domains nor its own eliminated domains.
Therefore \(S\eta\cup[\eta]\) is solved. On old solved
domains its substitution is \(A\eta\); on new \(\eta\)-domains
A was identity and the same composite gives the new image;
elsewhere both are identity. Hence the equality of substitutions
is global. In N, the sequence of helper eliminations also applies
their composite \(\eta\) to the whole skeleton, as claimed.
In B it applies none of them to P.

If \(\rho\) literally solves old S and Q, the old solved-form
identity gives \(A\rho=\rho\), and the helper factor identity
gives \(\eta\rho=\rho\). Thus
\(A\eta\rho=A\rho=\rho\). It also solves new S and
\(C\eta\) exactly when it solves old C. Conversely a literal
new-S solution satisfies \(\eta\rho=\rho\), since the
\([\eta]\) bindings are retained. Solving \(S\eta\) then
implies solving old S. Since Q is unified by \(\eta\),
\(Q\rho=(Q\eta)\rho\) is literal identity; and the same
calculation restores old C from \(C\eta\). The possible
deleted literal identities affect no solution. This proves exact
equivalence of the stated local literal conditions and supplies
the factor, not just existence of an unspecified MGU.

For Solve, Q=e A and every new literal witness has
\(A\rho=\rho\), so e is literally solved and hence equationally
solved. For Narrow, Q gives \(l\rho=u\rho\). The chosen
identity gives \(l\rho=_Rr\rho\), and congruence restores
the old equation e[u] from a solved new e[r]. In N the new
equation was also instantiated by \(\eta\); its witness's
\(\eta\rho=\rho\) restores the uninstantiated e[r]. The
other problem equations are restored in the same way. Thus every
successful block is backwards sound. A U failure means Q has no
literal solution; it does not exclude a different LP choice or
an equational solution that does not satisfy this chosen Q.
Termination is solely the finite U termination proof, not a
termination argument for repeated Solve/Narrow choices.
\end{proof}
\end{annexitem}

\begin{annexitem}{annex:p12.en.actual-lift}{thm:p12-en-actual-lift}{The chosen ground redex, fresh match, computed MGU and exact reduct}
\begin{proof}
Retain the finite witness on the externally fixed initial scope.
Choose the fresh rule variant also outside the finite support of
\(\rho\), which is possible in the countably infinite variable
supply and is permitted by the source's freshness requirement.
If the chosen redex occurred at a variable address of the basic
term, or strictly inside such a variable's inserted image, then
that image would contain an R-redex. A variable outside the
support is itself irreducible for an ordinary R; every relevant
image in the support is normal by the reduced-witness premise.
Therefore the redex address is nonvariable in the basic skeleton.

Write that subterm u. The actual ground reduction supplies a
matching map \(\gamma\) with \(l\gamma=u\rho\). Because l
is nonvariable, each variable occurring in l is matched to a
proper subterm of the innermost redex. Those subterms are normal.
If a variable repeats in l, the actual match gives the same
ground subterm at all occurrences. The ordinary-rule condition
\(\operatorname{Var}(r)\subseteq\operatorname{Var}(l)\)
ensures that every new right-side image is covered by this same
map, with no unassigned right-only variable. Let \(\rho'\)
extend \(\rho\) by these fresh normal images. It preserves all
old relevant values and remains reduced.

Freshness makes old S independent of those assignments. Hence
\(A\rho'=\rho'\) globally, and
\((l A)\rho'=l\gamma=u\rho=(u A)\rho'\).
Thus \(\rho'\) is an actual literal solution of Q and old C,S.
The actual U block cannot fail. Its factor identity yields
\(\eta\rho'=\rho'\), and the preceding composite theorem
gives \(A\eta\rho'=\rho'\). This is the explicit factor
through the computed solved substitution. No assumption of
``lifting completeness'' has entered the argument.

In B the new skeleton at that occurrence is e[r]. Its value
under \(\rho'\) is precisely the old ground context with
\(l\gamma\) replaced by \(r\gamma\), the originally chosen
ground reduct. In N the new skeleton is also instantiated by
\(\eta\), but \(\eta\rho'=\rho'\) makes its value the
same reduct. All other old instantiated terms retain their values.
The occurrence multiset of problem terms therefore decreases in
the source multiset ordering whenever that ground R-step decreases
its reduction ordering. Choosing the least indexed matching rule
is an actual finite choice and is preserved by the computed match;
it is not a choice by equational unifiability.

If the old witness R-solves P, the new ground reduct is R-equal
to its predecessor and preserves those equation solutions. The
same witness also solves the new constraints/bindings by the
literal computations. The trace compiler proves that the entire
block, its matching-first Decomposition and its answer factor
transport to N. All these assertions are local to the chosen
redex and witness; they do not manufacture the groundability
premise for every arbitrary parameter target.
\end{proof}
\end{annexitem}

\begin{annexitem}{annex:p12.en.terminal-normalization}{thm:p12-en-terminal-normalization}{Finite image reductions, idempotence and preservation of coverage}
\begin{proof}
Normalize the finitely many nonidentity images of A by the
assumed finite reduction sequences. Each rewrite is an equation
of R, so \(x A=_R x\widehat A\) on the domain, and equality
outside it is identity. Thus the two substitutions are pointwise
R-equal. By induction on terms, any input equation's A-instance
is R-equal to its \(\widehat A\)-instance. An answer's
equational soundness is consequently preserved.

In an ordinary rule, every right-side variable occurs on its left.
After matching that left side in a term, all variables of its
instantiated right side already occur in the matched subterm.
Therefore rewriting cannot introduce a variable absent from
the starting term. A's images contain no solved domain variable,
so neither do the normalized images. With the same domain and
identity elsewhere, \(\widehat A\) is idempotent.

For every finite factor \(\lambda\), substitution closure gives
\(x A\lambda=_R x\widehat A\lambda\).
Hence \(\theta\) is covered by A exactly when it is covered
by \(\widehat A\), on the same fixed original scope, with
the same factor. Applying this to each answer proves family
coverage in both directions. This argument requires only the
given image reductions. If R terminates on all those terms,
choosing redexes until none remain computes such forms; if it
is also confluent, all choices agree. Neither condition follows
merely from the existence of a successful narrowing trace.

No step asserted preservation of a literal active S-solution.
The displayed active-state example is reachable: start with
x=f(a),z=f(a), Constrain and eliminate x, then Constrain the
remaining equation. In that state x=z=f(a) literally solves
the old C,S, while x=a fails its old x-value under the new
binding. Equation f(a)=a holds only modulo R, not as finite-tree
identity. Terminal answer normalization concerns an equational
representative after C is empty, and is not the N elimination
operation or permission to normalize an active accumulator.
\end{proof}
\end{annexitem}

\begin{annexitem}{annex:p12.en.infinite-search}{thm:p12-en-infinite-search}{A convergent system with infinitely repeated legal fresh-skeleton choices}
\begin{proof}
Interpret both constants as1 and both unary operations as
\(v\mapsto v+1\) over positive integers. Then
\(f(s(x))\) has value \(x+2\), strictly greater than
\(f(x)\)'s \(x+1\); f(0) has value2, strictly greater
than a's1. Each rule decreases a positive integer term value
in every context, proving termination for finite terms.

Define a normal-form operation recursively. Constants and variables
are unchanged; normalize s(t) as s of the normalized child.
For f(t), first normalize its child to v, write it uniquely as
\(s^m(w)\) where w is not rooted at s, and return a if w=0,
otherwise f(w). The child w is normal; the returned term has no
redex at its root or below it. Both rule instances preserve this
operation: stripping the extra initial s in f(s(t)) gives the
same result as f(t), and f(0) and a both yield a. Recursion
shows that equal normalized children give equal normalized enclosing
contexts, so every contextual rewrite preserves the operation.
Every irreducible term is unchanged by it. Termination gives an
irreducible endpoint for every reduction, and preservation makes
all such endpoints the same. This proves confluence on all
finite terms, not just the original closed ground terms.

Let \(v_0=y\) and let \(v_1,v_2,\ldots\) be fresh variables.
After k choices of the first rule, the actual basic skeleton is
\(f(v_k)=^{?}a\), C is empty, and S maps
\(v_j\mapsto s^{k-j}(v_k)\) for \(0\leq j<k\).
At its next root LP use the fresh rule
\(f(s(v_{k+1}))\to f(v_{k+1})\).
The matching constraint first decomposes to
\(s(v_{k+1})=^{?}v_k\); Orient and elimination compute
\(\eta=\{v_k\mapsto s(v_{k+1})\}\).
The new skeleton is f(v_{k+1})=a, and source-order accumulator
composition gives exactly the same formula with k+1. The occurs
check passes and every root/nonvariable/freshness restriction holds.
The replaced f(v_{k+1}) is a newly introduced skeleton, not a
position inside an older solved image. In N, this new skeleton
contains no \(v_k\), so its \(\eta\)-instance is the same
term; alternatively the trace compiler gives the same sequence.

At any finite k choose the second rule instead. The match
f(0)=f(v_k) first decomposes to0=v_k, yielding
\(v_k\mapsto0\). The replacement equation is a=a; Solve
and Trivial finish with the returned y-image \(s^k(0)\).
Choosing the first rule at every k yields an infinite legal B
and N branch. There is no finite reduced ground witness for
that whole infinite branch, so it does not contradict the
source witness-descent proof. These concrete traces distinguish
termination of each finite syntactic block and of ground rewriting
from termination of the full narrowing search.

For completeness of this example's family, let t be any finite
image of y. The constructed normal-form operation is invariant
under each rule instance in either direction and under contexts,
so equational equality also preserves it. Conversely every term
reduces to its computed form by termination and the proved endpoint
uniqueness. Therefore \(f(t)=_Ra\) exactly when its computed
normal form is a. By the explicit f-case this is exactly when
the computed normal form of t is \(s^k(0)\) for some finite k.
It follows that every target unifier is R-equal on y to the
corresponding ground answer, which therefore covers it. Distinct
\(s^k(0)\) are distinct irreducible terms and hence not R-equal.
A factor substitution cannot change a ground answer. Thus different
members are incomparable, proving the infinite minimal complete
family on the initial scope. Discarding the source traces' auxiliary
bindings does not change this initial-scope comparison. In particular
no finite complete family exists for this problem: each of its
members, being itself a unifier, is R-equal on y to one ground
answer and cannot cover a different one: substituting in its
R-equality to that ground answer leaves the ground side fixed.
\end{proof}
\end{annexitem}
